%% kapa: exam=1 round=19 year=2008 period=1 form=A subject=civil_law subject_ko=민법 seq=16 number=16
%% source: Q-Net 감정평가사 기출문제 (cst00309) artlSeq=1001962 file=extracted/1008/1001962/2008년 제19회 감정평가사 1차시험 기출문제/1교시.hwp
%% conversion: kordoc 4.24.51 -> tools/gen1.py (auto); review=proofread-sonnet; flags=none
\begin{question}{16}
법률행위의 무효에 관한 설명으로 틀린 것은? (다툼이 있는 경우에는 판례에 의함)
\begin{choices}
  \item 변호사 아닌 자가 승소를 조건으로 그 대가로 소송당사자로부터 소송목적물 일부를 양도받기로 한 약정은 반사회적 법률행위로서 무효이다.
  \item 양도소득세 회피를 위하여 매매계약을 체결하였더라도 계약의 효력은 인정된다.
  \item 이미 매도된 부동산에 관하여 매도인의 채권자가 매도인의 배임행위에 적극 가담하여 저당권설정계약을 체결한 경우에 저당권설정계약은 유효이다.
  \item 조건이 선량한 풍속 기타 사회질서에 위반한 경우에는 조건뿐만 아니라 법률행위 전부가 무효로 된다.
  \item 토지거래허가 전의 거래계약이 정지조건부 계약이었는데 그 정지조건이 토지거래허가를 받기 전에 이미 불성취로 확정된 경우에는 허가 전의 거래계약관계는 확정적으로 무효로 된다.
\end{choices}
\end{question}
\answer{3}
